% (This file is included by thesis.tex; you do not latex it by itself.)

\newcommand{\newabstract}[1]{%
  \par\bigskip
  \csname otherlanguage*\endcsname{#1}%
  \csname captions#1\endcsname
  \item[\hskip\labelsep\scshape\abstractname.]
}

\begin{abstract}

% The text of the abstract goes here.  If you need to use a \section
% command you will need to use \section*, \subsection*, etc. so that
% you don't get any numbering.  You probably won't be using any of
% these commands in the abstract anyway.



In this work, we prove the Weighted Fundamental Lemma for Lie algebras for general (quasi-split) unramified $G$. The strategy is to adapt the proof of Chaudouard--Laumon from the split case. In particular, we define the notion of a $\xi$-stable Higgs bundle for a quasi-split reductive group scheme $G$ on a curve, and prove that the stack of such form a proper, finite-type Deligne--Mumford stack over the Hitchin base. We then adapt the global cohomological theorem of Chaudouard--Laumon to the quasi-split setting, and perform the required Frobenius-point count to obtain an identity of adelic orbital integrals. A local-to-global argument deduces the WFL for Lie algebras. In the last chapter, we perform a parallel reduction to deduce Waldspurger's nonstandard WFL from a global version which follows from the same geometric theory. Our work implies the Twisted Weighted Fundamental Lemma in good characteristic, which completes Moeglin--Waldspurger's stabilization of the twisted Arthur--Selberg trace formula. We therefore render unconditional the endoscopic classification for classical groups and an array of downstream consequences in automorphic theory.

\newabstract{french}
Dans ce travail, nous démontrons le lemme fondamental pondéré pour les algèbres de Lie pour un groupe $G$ quelconque. Notre stratégie consiste à adapter au cas général la preuve de Chaudouard--Laumon dans le cas déployé. Nous définissons en particulier la notion de fibré de Higgs $\xi$-stable pour un schéma en groupes réductif quasi-déployé $G$ sur une courbe, et démontrons que le champ de tels fibrés est un champ de Deligne--Mumford propre et de type fini sur la base de Hitchin. Nous adaptons ensuite au cadre quasi-déployé le théorème cohomologique global de Chaudouard--Laumon, puis effectuons le comptage de points fixes de Frobenius nécessaire afin d’obtenir une identité entre intégrales orbitales adéliques. Un argument local-global permet alors d’en déduire le lemme fondamental pondéré pour les algèbres de Lie. Dans le dernier chapitre, nous effectuons une réduction parallèle afin de déduire le lemme fondamental pondéré non standard de Waldspurger d’une version globale, laquelle découle de la même théorie géométrique. Notre travail implique le lemme fondamental pondéré tordu en bonne caractéristique, achevant ainsi la stabilisation, due à Mœglin--Waldspurger, de la formule des traces tordue d’Arthur--Selberg. Nous rendons par conséquent inconditionnelles la classification endoscopique des groupes classiques ainsi qu’une série de conséquences qui en découlent en théorie automorphe.

\end{abstract}
